\subsection{GP-1988-20: a finite-model audit of the three-dimensional Ising programme}
\label{dev:ising-audit}

\paragraph{Source and reconstruction.}
The closing instruction is
\ru{Дальше надо реш. задачу в непр. пределе.}
The notebook specifies neither a topology of convergence nor field
normalization. \ours\ A precise stronger target is the following.
Let $\mu_c$ be the free-boundary infinite-volume Gibbs measure of the
zero-field nearest-neighbour ferromagnetic Ising model on $\mathbb Z^3$
at its critical inverse temperature. Seek $c>0$ and
$0<\Delta_\sigma<3/2$ such that
\[
 \Phi_a(f)=c\,a^{3-\Delta_\sigma}
     \sum_{x\in\mathbb Z^3}\sigma_x f(ax),
 \qquad f\in C_c^\infty(\mathbb R^3),
\]
converges in law as a random distribution along the full limit
$a\downarrow0$ to a nonzero non-Gaussian field with conformally covariant
noncoincident correlations. A corresponding centred energy field,
its dimension $\Delta_\varepsilon$, and thermal relations require
additional assertions. This is an editorial conjectural target, not
a theorem proved here or a verbatim statement of the notes.

\paragraph{The checked literature.}
The Dotsenko--Polyakov chapter was bibliographically verified, but its
full proof and coefficient conventions were not obtained
\citep{dotsenkopolyakov1988fermion}.
In two dimensions, \citet[Theorem 1.2]{chelkak2015spin} prove convergence
of appropriately normalized critical square-grid spin correlations
with plus boundary conditions in simply connected domains.
Normalization, separated points and domain hypotheses matter.
This does not establish the three-dimensional target.

\citet[section 8.4]{duminilcopin2022hundred} explicitly distinguish
bootstrap information from proving a lattice scaling limit.
\citet[section 4]{kulske2025ising} supplies a later dated rigorous-status
discussion. The September 2026 effective-string paper concerns a gapped
bulk and low-energy interactions
\citep[section 2]{lopesmatospenedones2026strings}.
The January 2026 certified-exponent claim was identified, but its
complete certification was not independently audited
\citep{kai2026certified}. Even exponent certification requires a separate
argument to establish field convergence. No checked primary proof of
the target was found as of November 2026. This is a bounded search
conclusion, not a proof of universal absence.

\paragraph{Finite conventions.}
Let $G=(V,E)$ be a finite loopless undirected graph, each edge counted
once, with real temperature-independent couplings:
\[
 H(\sigma)=-\sum_{e=\{u,v\}}J_e\sigma_u\sigma_v,
 \qquad
 Z_G(\beta)=\sum_{\sigma\in\{\pm1\}^V}e^{-\beta H(\sigma)}.
\]
The notebook's simultaneous negative signs in the Hamiltonian and
bond-sum exponent conflict with this ferromagnetic convention.
Counting both edge orientations without one-half also changes the
normalization.

\begin{proposition}[Exact graph expansions]
\ours\ Put $t_e=\tanh(\beta J_e)$, and let $\partial A$ denote the
odd-degree vertices of $A\subseteq E$. Then
\[
 Z_G=2^{|V|}\prod_{e\in E}\cosh(\beta J_e)
       \sum_{\partial A=\varnothing}\prod_{e\in A}t_e.
\]
For distinct $u,v$,
\[
 \langle\sigma_u\sigma_v\rangle=
 \frac{\sum_{\partial A=\{u,v\}}\prod_{e\in A}t_e}
      {\sum_{\partial A=\varnothing}\prod_{e\in A}t_e}.
\]
For unit couplings, $q=e^{-2\beta}$ and $c(G)$ connected components,
\[
 Z_G=2^{c(G)}e^{\beta|E|}
       \sum_{D\in\mathcal B(G)}q^{|D|},
 \qquad
 \mathcal B(G)=\{\delta S:S\subseteq V\},
\]
where distinct cuts, rather than repeated vertex-set presentations,
are summed.
\end{proposition}

\begin{proof}
Expand each bond factor as
$\cosh(\beta J_e)(1+t_e\sigma_u\sigma_v)$.
The spin sum vanishes unless all vertex exponents are even; then it
equals $2^{|V|}$. Inserting $\sigma_u\sigma_v$ changes exactly the two
specified parity conditions. The denominator remains positive for
signed real couplings, since $Z_G$ and the cosh factors are positive.

For the last formula, a configuration has energy
$-|E|+2|\delta S|$. Two configurations give the same cut exactly when
their pointwise ratio is constant on each connected component.
Each distinct cut has $2^{c(G)}$ preimages, including freedom at
isolated vertices.
\end{proof}

For the four-cycle,
\[
 Z_{C_4}=2e^{4\beta}(1+6q^2+q^4)
       =16\cosh^4\beta(1+\tanh^4\beta).
\]
At $\beta=0$ this is sixteen, whereas the bare distinct-cut sum is eight.
The missing normalization is therefore detectable on four vertices.

\paragraph{Boundary conditions.}
For finite $\Lambda\subset\mathbb Z^3$ with plus exterior spins, let
$E_\Lambda$ contain every bond having at least one endpoint in $\Lambda$,
once. Then
\[
 e^{-\beta|E_\Lambda|}Z_\Lambda^+
 =\sum_{S\subseteq\Lambda}e^{-2\beta|\delta S|}.
\]
Dual plaquettes form cellular boundaries. They may meet along edges and
need not be embedded smooth surfaces. Free boundary conditions change
the admissible interface description.

\begin{proposition}[Local closure is insufficient on a torus]
\ours\ On $(\mathbb Z/L\mathbb Z)^3$, $L\ge3$, let $D$ consist of bonds
from $(L-1,y,z)$ to $(0,y,z)$. Every elementary plaquette meets $D$
evenly, but $D$ is not a disagreement set of periodic spins.
\end{proposition}

\begin{proof}
A plaquette crosses this seam zero or twice. An $x$-direction winding
cycle crosses it once. Every cut intersects every cycle evenly,
because crossing a cut changes side and a closed walk returns to its
initial side. Thus the seam is locally closed but is not a cut.
Three independent binary winding fluxes give eight cohomology sectors;
periodic spin disagreements occupy the trivial sector.
\end{proof}

\begin{proposition}[Connected energy covariance]
\ours\ Let $M=-H$ and $F=\log Z_G$. Then
\[
 F'=\langle M\rangle,\qquad F''=\operatorname{Var}(M)
 =\sum_{e,f}J_eJ_f
  \bigl(\langle s_es_f\rangle-\langle s_e\rangle\langle s_f\rangle\bigr),
 \quad s_{\{u,v\}}=\sigma_u\sigma_v.
\]
With Boltzmann's constant one, heat capacity per vertex is
$c_G=\beta^2F''/|V|$.

For translation-invariant couplings and the translation-invariant
finite-torus Gibbs measure, assign positive-direction bonds to their
initial vertex, so $M=\sum_x\varepsilon_x$. Then
\[
 |V|^{-1}F''=\sum_x\operatorname{Cov}(\varepsilon_0,\varepsilon_x).
\]
\end{proposition}

\begin{proof}
Differentiate the finite partition sum and its logarithm. The quotient
rule produces centring. Since $d\beta/dT=-\beta^2$ and $H=-M$,
differentiating mean energy gives the heat capacity.
Stationarity of both couplings and measure makes each displacement
occur $|V|$ times in the double covariance sum.
\end{proof}

A raw two-point function cannot replace connected covariance without
an explicit centring convention. An unsummed distance power cannot
equal a spatial sum while retaining its free spatial variable.

\begin{lemma}[Lattice shell growth]
\ours\ For integer $d\ge1$, $\Delta\ge0$ and integer $R\ge2$, put
\[
 S_R=\sum_{1\le\|x\|_\infty\le R}\|x\|_\infty^{-2\Delta}.
\]
If $2\Delta<d$, then $S_R\asymp R^{d-2\Delta}$.
If $2\Delta=d$, then $S_R\asymp\log R$.
If $2\Delta>d$, the sum converges to a finite positive limit.
\end{lemma}

\begin{proof}
The shell at radius $r$ contains
$(2r+1)^d-(2r-1)^d$ points. This polynomial has positive coefficients
and leading term $d2^d r^{d-1}$; consequently it is bounded above and
below by positive constant multiples of that power for $r\ge1$.
Comparison with $\sum_{r\le R}r^{d-1-2\Delta}$ gives all three cases.
For $d=3$ and $2\Delta=3$ the exact expression is
$24H_R+2H_R^{(3)}$.
\end{proof}

If an energy covariance has the proposed power law through correlation
length $\xi$, with adequate tail and short-distance control, and
$\xi\asymp|\beta-\beta_c|^{-\nu}$, the divergent case suggests
$\alpha=\nu(d-2\Delta_\varepsilon)$.
The additional assumption $\nu^{-1}=d-\Delta_\varepsilon$ gives the
notebook's ratio. These are conditional implications.
The marginal case permits logarithms; above it regular backgrounds
matter. None follows merely from writing a quartic Grassmann action.

\paragraph{Executed checks and independent review.}
\oursComputed\ Before the outage, exact enumeration checked all 1,099
labelled simple graphs with one to five vertices, unequal signed
rational couplings and two-point numerators on three graphs, and twelve
covariance examples at 70-digit precision. The cube has even-subgraph
polynomial
\[
 1+6t^4+16t^6+9t^8
\]
and distinct-cut polynomial
\[
 1+8q^3+15q^4+24q^5+32q^6+24q^7+15q^8+8q^9+q^{12}.
\]
A separately written JavaScript recovery program actually reran the
1,099 graph comparisons, cube polynomials, torus calculations and
333 hyperdeterminant degree comparisons after filesystem execution
became unavailable. Torus sides three and four have cut ranks 26 and 63,
and locally closed dimensions 29 and 66. These are finite checks,
not critical estimates.

An independent reader checked the parity, cut-fibre, covariance and
shell proofs. A separate hand calculation for a triangle disjoint from
an edge and an isolated vertex gives
\[
 Z=8e^{4\beta}(1+q+3q^2+3q^3)
  =64\cosh^4\beta(1+\tanh^3\beta),
\]
confirming the component factor. The reviewer requested the explicit
stationarity qualification now included. Proposed further computations
on six-vertex graphs and larger tori failed before process creation and
are not counted as executed.

The development supplies a reproducible audit of finite identities and
of the gap in the continuum argument. It proves no uniform critical
correlation bounds, tightness, uniqueness of a scaling limit or conformal
covariance. Recovering the exact published fermionic coefficients and
testing finite boxes with fixed boundaries and winding sectors would
be a further bounded task; a critical-limit proof would still remain
separate.
